\documentclass[11pt]{article}
\usepackage[T1]{fontenc}
\usepackage[margin=1in]{geometry}
\usepackage{amsmath,amssymb,amsthm}
\usepackage[hidelinks]{hyperref}
\newtheorem{theorem}{Theorem}
\newtheorem{lemma}{Lemma}
\title{Vertex-Disjoint Paths at optimal width on subcubic bipartite graphs}
\author{Atlas\thanks{The proof in this note was found using ChatGPT.}}
\date{}
\begin{document}
\maketitle
\begin{abstract}
Assuming deterministic \(\mathrm{FPT}\ne\mathrm{W[1]}\), we exclude effectively
\(2^{o(w\log w)}\) time, up to a polynomial input factor, for undirected
Vertex-Disjoint Paths parameterized by treewidth or by the width of a supplied
path decomposition. The exclusion holds on connected bipartite graphs of
maximum degree three with distinct leaf terminals and an optimal width
certificate.
\end{abstract}

\section{Statement}

In Vertex-Disjoint Paths, the input contains a finite simple undirected graph
and an explicit list of demand pairs of vertices. We ask for a path between
the endpoints of each pair, with the paths pairwise disjoint, including their
endpoints. The number of pairs is part of the input. Let \(B\) be the complete
binary input length, including the demand list and any supplied certificate.
For a nonnegative integer \(w\), put
\(\phi(w)=\max\{1,w\log_2(\max\{2,w\})\}\).

\begin{theorem}\label{vdp-width-w1}\hypertarget{vdp-width-w1}{}
Assume deterministic \(\mathrm{FPT}\ne\mathrm{W[1]}\).
For every fixed \(C\ge 0\) and finite nonnegative \(h\) admitting a computable
function \(T\) with \(h(w)\le \phi(w)/j\) whenever \(w\ge T(j)\) and \(j\)
is a positive integer, there is no correct deterministic solver running in
\(O(2^{h(w)}B^C)\) time for Vertex-Disjoint Paths. This holds separately
when \(w\) is intrinsic treewidth and when \(w\) is the width of a supplied path
decomposition. The hard instances are connected bipartite graphs of maximum
degree three, with distinct leaf terminals and equal treewidth and pathwidth.
The supplied path decomposition, when included, is optimal.
\end{theorem}

We construct a deterministic family of queries from a compact Boolean encoding
of Clique. Small width follows by grouping the variables before constructing
the graph. We then simulate directed paths and pad the graph to certify the
width.

\section{From Clique to permutation hitting}

We use the \(\mathrm{W[1]}\)-hardness of Clique parameterized by the requested
size \(k\)~\cite{df}. For a graph with \(n\ge 2\) vertices, use \(k\) blocks of
\(\lceil\log_2 n\rceil\) Boolean variables to encode the chosen vertices.
Forbid each unused code in each block, and for each pair of blocks forbid every
pair of valid codes representing equal or nonadjacent vertices. Each prohibition
is a clause on the affected blocks, falsified exactly by the forbidden pattern.
The formula is satisfiable exactly when the graph has a \(k\)-clique. It has
\(m=k\lceil\log_2 n\rceil\) variables, and its complete length and construction
time are both bounded by \(F(k)n^a\) for computable \(F\) and a fixed \(a\).
We decide trivial Clique inputs directly.

Take an explicit CNF over \(m\) listed variables. Decide instances with fewer
than \(16\) variables or an empty clause directly, and likewise decide a
formula with no clauses. Put
\[
 b=\lfloor (\log_2 m)/2\rfloor,\qquad
 q=2^b,\qquad s=\lceil m/b\rceil,\qquad d=\max\{s,q\}.
\]
Pad the variable list to \(sb\) and partition it into \(s\) groups of size
\(b\). In row \(i\) of a \(d\times d\) table, the first \(q\) columns encode
assignments to group \(i\), for each \(i\le s\). Require a row selection to
hit the set of these \(q\) cells for each real row. For each clause, require
it to hit the set of all real cells whose group assignment satisfies a
literal of the clause. Thus a mapping selecting one column per row hits all
sets exactly when its real rows encode a satisfying assignment. Dummy
rows have no domain constraint. The set family is explicit and has polynomial
size in the formula length and \(m\). We have \(d=O(m/\log m)\).

Use the deterministic constructive coloring family of Lokshtanov, Marx, and
Saurabh~\cite[Lemma 2.8 and its proof]{lms}. It enumerates maps
\(f\colon[d]\times[d]\to[d+1]\) in time
\(2^{O(d\log\log(\max\{d,4\}))}d^{O(1)}\). For any mapping selecting one cell
per row, some \(f\) gives the selected cells distinct colors and gives each
selected cell a color absent from the other cells of its row. Indeed, the
graph formed by a clique on the selected cells and the stars joining each
selection to its row is the cactus-grid graph of that lemma.

For each \(f\), retain the cells whose colors are unique within their rows, and
replace the column of every retained cell by its color. Apply this operation to
every set and additionally require a hit on the retained cells of each row.
Add one unconstrained row to make a square of size \(D=d+1\).
For any permutation hitting these new sets, the row-color choice in each
original row has a unique preimage, which gives an original selection hitting
all old sets. Conversely, for a selection hitting all old sets, the map
\(f\) guaranteed by the coloring family retains every chosen cell and assigns
them distinct colors. Assign the unused color to the new dummy row. This
gives a permutation hitting all new sets. The number of queries and the time
to enumerate the family are \(2^{o(m)}\) up to a fixed polynomial factor in
the explicit formula length, with a computable sublinear exponent bound
\(O(m\log\log m/\log m)\).

\section{Testing the same permutation in every set}

We give the directed construction of Lokshtanov, Marx, and Saurabh
\cite[Section 5]{lms} on the permutation interface just obtained. For a
set \(S\subseteq [D]\times[D]\), make vertices \(a_i,b_j,u_{i,j}\) and
arcs \(a_i\to u_{i,j}\to b_j\). For each row \(i\), make a directed path
\[
 P_i=(c_{i,0},d_{i,1},w_{i,1},c_{i,1},\dots,d_{i,D},w_{i,D},c_{i,D}).
\]
For each cell \((i,j)\), add vertices \(f_{i,j},\ell_{i,j},r_{i,j}\)
and arcs
\[
 b_j\to f_{i,j}\to c_{i,j},\qquad
 \ell_{i,j}\to f_{i,j}\to r_{i,j},\qquad
 \ell_{i,j}\to c_{i,0},\qquad c_{i,j-1}\to r_{i,j}.
\]
Make a source \(s\) and a sink \(t\), and add \(s\to d_{i,j}\to t\)
for \((i,j)\in S\). The demands are \((a_i,c_{i,D})\) for each row,
\((\ell_{i,j},r_{i,j})\) for each cell, and \((s,t)\). All demand endpoints
are distinct. Each source has indegree zero and each sink has outdegree zero.

A row path choosing \(u_{i,j}\) must follow \(b_j,f_{i,j}\) into its own
row and then traverse the suffix from \(c_{i,j}\) to \(c_{i,D}\). Any other
route from \(b_j\) would enter a different row or end at another demand's
terminal. Distinct row paths use distinct vertices \(b_j\). Their choices
therefore define a permutation \(\pi\).
The local demand for \((i,\pi(i))\) must follow the prefix through
\(c_{i,\pi(i)-1}\) since its \(f\)-vertex is occupied. The prefix and suffix
together occupy all \(d_{i,j}\) and \(w_{i,j}\) except those at
\(j=\pi(i)\). The test path must start through an unoccupied \(d_{i,j}\),
so a hit exists. Conversely, route the rows and the special local demands
as above, and route all other local demands directly through their own
\(f\)-vertices. The test path uses \(s\to d_{i,\pi(i)}\to t\) for a hit.
These routes are disjoint and leave exactly \(w_{i,\pi(i)}\) free among
the output vertices of each row.

For the full set family, make one such graph per set. Identify each
\(w_{i,j}\) of a graph with \(u_{i,j}\) of the next graph. Retain all demands.
The routes for a common permutation combine since the only input used in
each row of the next graph is precisely the free output of the preceding
one. For the converse, induct by gadget index to show that paths for its
demands stay inside it, where shared vertices belong to both adjacent
gadgets. When a path follows an arc from a shared vertex into the next
gadget, it reaches a \(b_j\) there. To return it would have to reach an
input of that gadget. Within the gadget only indegree-zero row sources have
arcs into its inputs. An excursion further forward cannot change this:
each newly entered gadget has the same interface and returning first from
the furthest gadget reached would require reaching one of its inputs
through such a source.
For a backward excursion into the preceding gadget, the prefix and suffix
demands of the normal solution occupy all but at most one output per row.
A directed path through that previous gadget can return from that
row only along the same row path. Returning through the sole available
output would repeat a vertex. The first gadget has no preceding gadget.
This proves the induction. The normal routes in two successive gadgets
choose the same column per row, since the used input of the second must
be the available output of the first. Thus a common permutation hits all
sets.

We supply a path decomposition of the underlying undirected graph. Initially
process the first inputs \(u_{i,j}\) one at a time with all \(a_i,b_j\) of
the first gadget. In a gadget block retain its \(a_i,b_j,s,t\) and the
next gadget's \(a_i,b_j\), if it exists. Process the rows in order, retaining
the full path \(P_i\) throughout its row block. For each column add the
vertices \(f_{i,j},\ell_{i,j},r_{i,j}\) in one bag. All internal edges are
covered, and the shared output vertices meet the next gadget's \(a_i,b_j\)
in the pertinent row block. The next gadget's \(a_i,b_j\) persist across
the two gadget blocks. The row paths and local vertices occur on consecutive intervals by
construction. The largest bag has size at most \(7D+6\). The graph and the
listed bags have polynomial size in the table and the set list.

\section{Simulating direction on a subcubic graph}

Let \(w_0\) be the width of the directed instance's path decomposition.
We first modify the digraph so every vertex has indegree and outdegree at most
two. Assign each arc to a bag containing both endpoints and order the arcs by
bag and a fixed tie order. Replace each vertex by a center with a directed
chain of incoming ports leading to it and a directed chain of outgoing ports
leading away from it. Order ports in each chain by their assigned arcs.
Every original arc now connects its outgoing port at its tail to its incoming
port at its head. Put the demands on the corresponding centers. A traversal
from an incoming incidence to an outgoing one uses the center, so directed
routes expand and contract to give equivalent solutions. Any route that
starts or ends at a demanded center also uses that center. A source center
has indegree zero and at most one outgoing arc. A sink center has outdegree
zero and at most one incoming arc.

For the width, keep every center throughout its original bag interval.
Process the assigned arcs at their bags. For each live center and each port
chain retain the last introduced port until the next port of that chain
arrives, or until the final incident chain edge is covered. Introduce the
two fresh ports for an arc together to cover its replacement edge and the
chain edges to retained predecessors. The first outgoing port and the last
incoming port meet the center. Each bag contains at most three retained
vertices per old live vertex and two new ports. Every port occurrence is
consecutive. This gives a polynomial-length decomposition of width at most
\(3(w_0+1)+1\).

For each split digraph vertex \(v\), make adjacent vertices \(v^-\) and
\(v^+\). Replace an arc \(u\to v\) by an undirected path
\(u^+,p_{uv},q_{uv},v^-\) with private internal vertices. A main demand
from a source \(s\) to a sink \(t\) now joins \(s^-\) to \(t^+\), both
leaves. If \(v\) has two incoming arcs with private pairs \(p_i,q_i\) for
\(i=1,2\), add \(a\) adjacent to both \(p_i\) and \(b\) adjacent to both
\(q_i\). Put a private leaf \(a'\) at \(a\) and \(b'\) at \(b\), and add
the auxiliary demand \((a',b')\). All vertices have degree at most three,
and all demand endpoints are distinct leaves.

A directed routing extends to main routes using the natural arc paths.
At most one incoming arc path at each \(v\) is in use, so the auxiliary
route can use \(a',a,p_i,q_i,b,b'\) on a free arc path. Conversely, an
auxiliary route uses \(a,b\) and at least one of the two \(p_i\). A main
route using any private pair must traverse the complete arc path because
the only branches lead to the occupied \(a,b\), and neither private vertex
is a terminal. At most one incoming arc path at each \(v\) is therefore
available to the main routes. A main route starts by crossing from its
source leaf to the source's positive copy. After an outgoing arc path
it reaches another negative copy and must cross to the positive copy,
since no other incoming arc path is available. It continues in this way
until its own sink leaf. The projected directed routes are disjoint since
each visited vertex uses both copies.

Double the width decomposition's bags and insert a local bag adding two
private vertices for each arc at a bag containing its endpoints. At an
indegree-two head, extend the earlier private pair through the local bag
of the later pair. The extension stays inside the interval of the negative
copy of the common head, which is present in both local bags. Each head
requires at most two extended vertices per bag. Where the two pairs meet,
insert copies adding \(a,a'\) and \(b,b'\) in separate bags. This covers
all auxiliary edges and preserves contiguity. The width is \(O(w_0+1)\)
and the listed bag size is polynomial in the preceding input.

\section{Certifying the width on a connected bipartite graph}

Subdivide each edge once, using local bags to add each subdivision vertex.
This preserves routing and the terminal leaves, and costs one in width.
It makes the graph bipartite. Discard isolated vertices. If a demand
crosses between components we reject the query directly. Otherwise choose
one subdivision vertex in each component and attach it by a bridge to
its own vertex on a new spine path. The attachments preserve bipartiteness
by choosing the coloring of each component independently, and preserve
maximum degree three. A simple path with its two endpoints in a component
must stay there because it has only one edge to the rest of the graph.
Process component decompositions in spine order, retaining the appropriate
spine vertex in every bag for its component and adding bags for the spine
edges. The width increases by at most one. If no component remains, decide
the query directly.

Let \(Q\) be the maximum of three and the actual width of this decomposition.
We construct a connected bipartite padding graph with both widths \(Q\).
Take \(Q+1\) disjoint paths indexed \(0,\ldots,Q\). For each pair \(i<j\),
make one port on path \(i\) and one on path \(j\) and join them by an edge.
Order the ports on each path by lexicographic order of the pairs \((i,j)\).
Color the port on the lower-index path black and the other one white.
Between consecutive ports use a path of length one for different colors
and two for equal colors. This gives a proper bipartition and maximum
degree three. Contracting the indexed paths yields a \(K_{Q+1}\) minor,
so treewidth is at least \(Q\).

For a path decomposition, process pairs in lexicographic order. Retain
the current vertex on each started path while it has another pair
to process. Advance the two relevant paths to their next ports one at
a time, introducing one vertex before forgetting its predecessor.
Once the ports meet, forget any final ports. Introduce an unstarted path
at its first port. During the pairs with first index zero, advance path
zero before introducing the other path. Before pair \((0,j)\) there
are at most \(j\) retained vertices, and at most \(j+1\) appear in any
bag for this pair. The first pair starts with the two ports. After
\((0,Q)\) forget the final port of path zero, so subsequent bags have
at most \(Q\) retained vertices plus one during an advance. All edges are
covered and every vertex has a consecutive lifetime, so the bag size
bound is \(Q+1\).

Attach the first port on path zero by a bridge to the last spine vertex.
Both endpoints have room for the added edge under the degree bound, and
we can flip the padding bipartition if needed. The padding has no terminals
and creates no new routes between old terminals. Concatenate the
decompositions with a bridge bag between the last spine component and
the first port bag. We obtain an optimal path decomposition of width \(Q\)
because of the clique minor. This is also an optimal tree decomposition.
All transformations have polynomial running time and output length in
the initial directed instance and its decomposition.

\section{Running time}

For each formula query, the constructed width satisfies
\(\phi(Q)\le A m+A\) for a fixed positive integer \(A\). All query bit
lengths are bounded by a fixed polynomial in the formula length and \(m\).
For a Clique instance, they are bounded by \(F_1(k)n^{a_1}\) for computable
\(F_1\) and fixed \(a_1\). The enumeration and construction overhead has
the same form: its effectively sublinear overhead in \(m\) is at most
\(2^{m/(2k)}\) beyond a computable cutoff depending on \(k\), and its
finite prefix has a computable bound.

Suppose a solver as in Theorem~\ref{vdp-width-w1} exists. Both applications
measure the same \(Q\) on our queries. For \(j=2A(k+1)\) and
\(Q\ge T(j)\), we have
\[
 h(Q)\le (A k\lceil\log_2 n\rceil+A)/j
       \le \log_2 n+1.
\]
For the exceptional widths, hardwire a single integer \(D_0\) bounding
the finite values of \(h(w)\) for \(w<T(1)\). Then
\(h(w)\le \phi(w)+D_0\) for all \(w\). Consequently the range
\(Q<T(j)\) has a computable upper bound on \(h(Q)\) as a function of \(k\),
even if \(h\) itself is uncomputable. The polynomial input factor has
a fixed exponent. Running the solver on every query thus gives a
deterministic FPT algorithm for Clique, a contradiction.

\begin{thebibliography}{9}
\bibitem{df} R. G. Downey and M. R. Fellows.
Fixed-parameter tractability and completeness II: On completeness for W[1].
\emph{Theoretical Computer Science} 141(1--2), 109--131, 1995.
\href{https://doi.org/10.1016/0304-3975(94)00097-3}{DOI}.
\bibitem{lms} D. Lokshtanov, D. Marx, and S. Saurabh.
Slightly Superexponential Parameterized Problems.
\emph{SIAM Journal on Computing} 47(3), 675--702, 2018.
\href{https://doi.org/10.1137/16M1104834}{DOI}.
\end{thebibliography}
\end{document}
